% =============================================================================
% Prakata (Kata Pengantar)
%
% PPTA 2026: Persantunan diberikan kepada pihak yang berkontribusi langsung,
% termasuk sumber dana bila ada.
% =============================================================================

\ipbprelimheading{PRAKATA}

Puji dan syukur penulis panjatkan kepada Tuhan Yang Maha Esa atas segala
karunia-Nya sehingga skripsi ini berhasil diselesaikan. Judul dalam penelitian
yang dilaksanakan sejak bulan \BulanMulaiPenelitian{} \TahunMulaiPenelitian{}
sampai bulan \BulanSelesaiPenelitian{} \TahunSelesaiPenelitian{} ini ialah
``\JudulSkripsi''.

Terima kasih penulis ucapkan kepada para pembimbing,
\NamaPembimbingSatu{} dan \NamaPembimbingDua, yang telah membimbing
dan banyak memberi saran. Ucapan terima kasih juga disampaikan kepada
pembimbing akademik, moderator seminar, dan penguji luar komisi pembimbing.
Penghargaan penulis sampaikan kepada \ldots{} (nama lengkap dan gelar dari
lembaga/instansi/perusahaan yang telah memberi izin penelitian), \ldots{}
(nama dan gelar atau bapak/ibu jika tidak ada gelar) beserta staf Laboratorium
\ldots{} yang telah membantu selama pengumpulan data. Ungkapan terima kasih
juga disampaikan kepada ayah, ibu, serta seluruh keluarga yang telah
memberikan dukungan, doa, dan kasih sayangnya \ldots{} dan seterusnya.

\ifdefempty{\SumberDana}{}{Penelitian ini didukung oleh \SumberDana.}

Semoga tugas akhir ini bermanfaat bagi pihak yang membutuhkan dan bagi
kemajuan ilmu pengetahuan.

\vspace{14pt}

\noindent\textcolor{red}{Catatan: Ucapan terima kasih diberikan hanya kepada
pihak-pihak yang secara langsung berkontribusi terhadap pengumpulan data dan
penulisan tugas akhir.}

\vspace{2\baselineskip}

\begin{flushright}
\KotaUniversitas, \BulanTahunDokumen\\[\baselineskip]
\NamaPenulis
\end{flushright}
